\documentclass[../Odyssey_Project.tex]{subfiles}

\begin{document}
\setcounter{section}{1}
\section{Automorphisms}

\subsubsection*{The Automorphism Group}
\begin{itemize}
    \item The set of automorphisms of a group $G$ is denoted $\operatorname{Aut}(G)$.
    \item If $\varphi, \rho \in \operatorname{Aut}(G)$, then $\varphi \circ \rho \in \operatorname{Aut}(G)$, so composition is a binary operation on $\operatorname{Aut}(G)$.
    \item Under composition, $\operatorname{Aut}(G)$ is a group: identity is the trivial automorphism, and the inverse of $\varphi$ is its set-map inverse $\varphi^{-1}$.
    \item Call $\operatorname{Aut}(G)$ with this operation the \emph{automorphism group} of $G$; we may write $\varphi\rho$ for $\varphi \circ \rho$.
\end{itemize}

\subsubsection*{Inner Automorphisms, $\operatorname{Inn}(G)$, and the Center}
\begin{itemize}
    \item Each $g \in G$ defines a conjugation homomorphism $\varphi_g : G \to G$ by $\varphi_g(x) = gxg^{-1}$.
    \item Each $\varphi_g$ is an automorphism: for $x \in G$, $x = \varphi_g(g^{-1}xg)$, and if $\varphi_g(x) = \varphi_g(y)$ then $x = y$ by cancellation. These are the \emph{inner automorphisms} of $G$.
    \item $\varphi_g \varphi_h = \varphi_{gh}$ for $g,h \in G$, since $g(hxh^{-1})g^{-1} = (gh)x(gh)^{-1}$. Hence there is a homomorphism $G \to \operatorname{Aut}(G)$ sending $g$ to $\varphi_g$.
    \item Its image is the \emph{inner automorphism group} of $G$, denoted $\operatorname{Inn}(G)$; its kernel is the \emph{center} of $G$, denoted $Z(G)$. Observe
    \[
    Z(G) = \{g \in G \mid \varphi_g(x) = x \text{ for all } x \in G\}
         = \{g \in G \mid gx = xg \text{ for all } x \in G\},
    \]
    so $Z(G)$ consists of those elements of $G$ which commute with every element of $G$.
    \item Clearly $G$ is abelian iff $Z(G) = G$.
\end{itemize}

\subsubsection*{Outer Automorphisms, $\operatorname{Out}(G)$}
\begin{itemize}
    \item If $\sigma \in \operatorname{Aut}(G)$ and $\varphi_g \in \operatorname{Inn}(G)$, then $\sigma \varphi_g \sigma^{-1} = \varphi_{\sigma(g)}$. Hence $\operatorname{Inn}(G) \trianglelefteq \operatorname{Aut}(G)$.
    \item The quotient $\operatorname{Aut}(G)/\operatorname{Inn}(G)$ is the \emph{outer automorphism group} of $G$, denoted $\operatorname{Out}(G)$.
    \item ``Outer automorphism'' refers to automorphisms of $G$ which are not inner and which thus have non-trivial image in $\operatorname{Out}(G)$.
    \item If $G$ is abelian, then all non-trivial automorphisms of $G$ are outer in this sense, since $\operatorname{Inn}(G) = 1$.
\end{itemize}

\subsubsection*{Automorphism Groups of Cyclic Groups}
\begin{itemize}
    \item Let $G = \langle x \rangle \cong \Z$, and let $\varphi$ be an automorphism of $G$. Then $\varphi(x)$ must generate $G$; the only generators are $x$ and $x^{-1}$. So $\varphi$ either fixes each element or sends each to its inverse, and $\operatorname{Aut}(G) \cong \Z_2$.
    \item Let $n \in \N$ and $G = \langle x \rangle \cong \Z_n$. Suppose $\varphi$ is an endomorphism of $G$. Then $\varphi(x) = x^m$ for some $0 \leq m < n$, so $\varphi$ sends every element to its $m$th power. Hence $G$ has exactly $n$ endomorphisms, namely the $m$th power maps $\sigma_m$ for $0 \leq m < n$.
\end{itemize}

\begin{proposition}
\label{prop:2.1}
Let $G = \langle x \rangle \cong \Z_n$ for $n \in \N$, and for each $0 \leq m < n$ let $\sigma_m$ be the endomorphism of $G$ sending $x$ to $x^m$. Then $\operatorname{Aut}(G)$ consists precisely of those $\sigma_m$ for which $m \neq 0$ and $\gcd(m,n) = 1$. Furthermore, $\operatorname{Aut}(G)$ is abelian and is isomorphic with the group $(\Z/n\Z)^{\times}$ of units of the ring $\Z/n\Z$.
\end{proposition}
\begin{proof}
    The map $\sigma_0$ has trivial image, so it is not an automorphism. Let $1 \leq m < n$ and consider $\sigma_m$. Suppose first that $\gcd(m,n) = 1$. Then there exist $a,b \in \Z$ such that $am + bn = 1$, and hence
    \[
        \sigma_m(x^a) = x^{am} = x^{1 - bn} = x(x^n)^{-b} = x.
    \]
    Thus $x \in \operatorname{im}\sigma_m$, so $\sigma_m$ is surjective. Since $G$ is finite, $\sigma_m$ is also injective, and therefore $\sigma_m \in \operatorname{Aut}(G)$. \\
    Conversely, suppose that $\sigma_m \in \operatorname{Aut}(G)$. Since $\sigma_m$ is surjective, there is some $a \in \Z$ such that $x = \sigma_m(x^a) = x^{am}$. Hence $x^{am - 1} = 1$, so $am - 1 = bn$ for some $b \in \Z$. It follows that any common divisor of $m$ and $n$ divides $1$, and therefore $\gcd(m,n) = 1$. This proves the first assertion. \\
    Now let $1 \leq m_1,m_2 < n$ with $\gcd(m_1,n) = \gcd(m_2,n) = 1$. If $t$ is chosen with $1 \leq t < n$ and $m_1m_2 \equiv t \pmod{n}$, then
    \[
        (\sigma_{m_1}\sigma_{m_2})(x) = \sigma_{m_1}(x^{m_2}) = x^{m_1m_2} = x^t.
    \]
    Hence $\sigma_{m_1}\sigma_{m_2} = \sigma_t = \sigma_{m_2}\sigma_{m_1}$, so $\operatorname{Aut}(G)$ is abelian. Finally, the map
    \[
        \operatorname{Aut}(G) \to (\Z/n\Z)^{\times}, \qquad \sigma_m \mapsto m + n\Z
    \]
    is a bijection by the first assertion, and the computation above shows that it preserves multiplication. Hence $\operatorname{Aut}(G) \cong (\Z/n\Z)^{\times}$.
\end{proof}

\subsubsection*{The Totient and Primitive Roots}
\begin{itemize}
    \item The \emph{totient} of $n \in \N$ is the number of positive integers less than $n$ and coprime to $n$ (the value at $n$ of the Euler phi-function).
    \item If $n = p_1^{a_1} \cdots p_r^{a_r}$ with the $p_i$ distinct primes, then the totient of $n$ is $(p_1^{a_1} - p_1^{a_1 - 1}) \cdots (p_r^{a_r} - p_r^{a_r - 1})$. To see this, note that $m$ is coprime to $n$ iff no $p_i$ divides $m$. Among $1,\dots,n$, exactly $n/p_i$ integers are divisible by $p_i$, exactly $n/(p_i p_j)$ are divisible by both $p_i$ and $p_j$, and so on. Hence inclusion-exclusion gives
    \[
        \varphi(n)
        = n \prod_{i = 1}^r \left(1 - \frac{1}{p_i}\right)
        = \prod_{i = 1}^r \left(p_i^{a_i} - p_i^{a_i - 1}\right).
    \]
    \item By Proposition~\ref{prop:2.1}, the order of $\operatorname{Aut}(\Z_n)$ is the totient of $n$. In particular, $|\operatorname{Aut}(\Z_p)| = p - 1$ when $p$ is prime.
\end{itemize}

\begin{proposition}
\label{prop:2.2}
Let $p$ be a prime. Then $\operatorname{Aut}(\Z_p) \cong \Z_{p-1}$.
\end{proposition}
\begin{proof}
    Let $F = \mathbb{F}_p$. By Proposition~\ref{prop:2.1}, $\operatorname{Aut}(\Z_p)$ is isomorphic with $F^{\times} = (\Z/p\Z)^{\times}$. Thus it suffices to show that $F^{\times}$ is cyclic of order $p - 1$. \\
    For each divisor $d$ of $p - 1$, define
    \[
        f_d = \lvert \{x \in F^{\times} \mid o(x) = d\} \rvert,
        \qquad
        z_d = \lvert \{x \in \Z_{p - 1} \mid o(x) = d\} \rvert.
    \]
    We first show that $f_d \leq z_d$ for every $d \mid (p - 1)$. If $x \in F^{\times}$ has order dividing $d$, then $x^d = 1$, so $x$ is a root of $T^d - 1 \in F[T]$. Since this polynomial has degree $d$, it has at most $d$ roots in $F$. Now suppose $F^{\times}$ has an element $x$ of order $d$. Then the $d$ elements of $\langle x \rangle$ are all roots of $T^d - 1$, so they are all the roots. Hence every element of $F^{\times}$ of order $d$ lies in $\langle x \rangle \cong \Z_d$. Therefore either $f_d = 0$, or $f_d$ is the number of elements of order $d$ in a cyclic group of order $d$. By Theorem~\ref{thm:1.4}, the cyclic group $\Z_{p - 1}$ has a unique subgroup of order $d$; every element of order $d$ generates a subgroup of order $d$, so all elements of order $d$ in $\Z_{p - 1}$ lie in this unique subgroup. This subgroup is cyclic of order $d$, and hence $f_d \leq z_d$. \\
    Since every element of $F^{\times}$ has order dividing $\lvert F^{\times} \rvert = p - 1$, and every element of $\Z_{p - 1}$ has order dividing $p - 1$, we have
    \[
        \sum_{d \mid (p - 1)} f_d = \lvert F^{\times} \rvert = p - 1
        = \lvert \Z_{p - 1} \rvert = \sum_{d \mid (p - 1)} z_d.
    \]
    Since $f_d \leq z_d$ for every $d \mid (p - 1)$ and the sums are equal, we must have $f_d = z_d$ for every $d \mid (p - 1)$. In particular, $f_{p - 1} = z_{p - 1}$. But $\Z_{p - 1}$ has an element of order $p - 1$, so $z_{p - 1} > 0$, and hence $f_{p - 1} > 0$. Therefore $F^{\times}$ has an element of order $p - 1$, so $F^{\times} \cong \Z_{p - 1}$. \\
    \(\therefore \operatorname{Aut}(\Z_p) \cong \Z_{p - 1}\).
\end{proof}

\subsubsection*{Primitive Roots Modulo $n$}
\begin{itemize}
    \item Let $G = \langle x \rangle \cong \Z_n$, $1 \leq m \leq n$, $\gcd(m,n) = 1$. Induction gives $(\sigma_m)^k(x) = x^{m^k}$.
    \item The order of $\sigma_m$ is the least positive $k$ such that $x^{m^k} = x$, i.e.\ the smallest $k \in \N$ with $m^k \equiv 1 \pmod{n}$.
    \item If the order of $\sigma_m$ equals the totient of $n$, we call $m$ a \emph{primitive root modulo $n$}.
    \item $\operatorname{Aut}(\Z_n)$ is cyclic iff there is a primitive root modulo $n$.
    \item For composite $n$, the structure of $\operatorname{Aut}(\Z_n)$ belongs to number theory more than group theory.
\end{itemize}

\begin{theorem}
\label{thm:2.3}
For $n > 1$, $\operatorname{Aut}(\Z_n)$ is cyclic iff $n = 2$ or $4$, or $n = p^k$ or $2p^k$ for some odd prime $p$ and some $k \in \N$.
\end{theorem}
\begin{proof}
    By Proposition~\ref{prop:2.1},
    \[
        \operatorname{Aut}(\Z_n) \cong (\Z/n\Z)^{\times}.
    \]
    Under this isomorphism, $\sigma_m$ corresponds to the unit $m + n\Z$, and the order of $\sigma_m$ is the least $k \in \N$ such that $m^k \equiv 1 \pmod{n}$. Hence $\operatorname{Aut}(\Z_n)$ is cyclic iff there is a primitive root modulo $n$. It remains to prove the number-theoretic classification of those $n > 1$ which have primitive roots. \\
    We use Euler's theorem in the following form: if $\gcd(a,q) = 1$, then $a^{\varphi(q)} \equiv 1 \pmod{q}$. First suppose $q = 2^s$ with $s \geq 3$. If $a$ is odd, then $a^2 \equiv 1 \pmod{8}$. We claim that
    \[
        a^{2^{s-2}} \equiv 1 \pmod{2^s}
    \]
    for all $s \geq 3$. The case $s = 3$ is the congruence just observed. If $a^{2^{s-2}} = 1 + b2^s$, then
    \[
        a^{2^{s-1}} = (1 + b2^s)^2 \equiv 1 \pmod{2^{s+1}},
    \]
    so the claim follows by induction. But $\varphi(2^s) = 2^{s-1}$, so every odd residue modulo $2^s$ has order at most $2^{s-2} < \varphi(2^s)$. Thus $2^s$ has no primitive root for $s \geq 3$. \\
    Next suppose $\gcd(u,v) = 1$ with $u,v > 2$. Let $a$ be coprime to $uv$, and put
    \[
        h = \operatorname{lcm}(\varphi(u), \varphi(v)).
    \]
    Since $u,v > 2$, both $\varphi(u)$ and $\varphi(v)$ are even (the reduced residue classes pair as $b$ and $-b$), so
    \[
        h \leq \frac{\varphi(u)\varphi(v)}{2} = \frac{\varphi(uv)}{2}.
    \]
    Euler's theorem gives $a^h \equiv 1 \pmod{u}$ and $a^h \equiv 1 \pmod{v}$; since $\gcd(u,v) = 1$, this implies $a^h \equiv 1 \pmod{uv}$. Hence no unit modulo $uv$ can have order $\varphi(uv)$, and $uv$ has no primitive root. \\
    These two obstructions leave only $2$, $4$, odd prime powers, and twice odd prime powers. Indeed, if
    \[
        n = 2^a p_1^{e_1} \cdots p_t^{e_t}
    \]
    is the prime factorization of $n > 1$, then two distinct odd prime factors are ruled out by the previous paragraph, as is $2^a p^e$ with $a \geq 2$ and $p$ odd. The pure powers $2^a$ with $a \geq 3$ were ruled out in the paragraph before that. Therefore any $n > 1$ with a primitive root must be one of $2$, $4$, $p^e$, or $2p^e$ with $p$ odd. \\
    It remains to see that all numbers in this list actually have primitive roots. Clearly $1$ is a primitive root modulo $2$, and $3$ is a primitive root modulo $4$. Now let $p$ be an odd prime. By Proposition~\ref{prop:2.2}, there is a primitive root $r$ modulo $p$. We may choose such an $r$ with
    \[
        r^{p-1} \not\equiv 1 \pmod{p^2}.
    \]
    For if a chosen primitive root $r$ satisfies $r^{p-1} \equiv 1 \pmod{p^2}$, then $r + p$ is still primitive modulo $p$, and
    \[
        (r+p)^{p-1} \equiv r^{p-1} + (p-1)pr^{p-2}
        \equiv 1 - pr^{p-2} \not\equiv 1 \pmod{p^2}.
    \]
    We claim that this $r$ is primitive modulo $p^e$ for every $e \in \N$. For $e = 1$ this is how $r$ was chosen. For $e \geq 2$, we first show by induction that
    \[
        r^{p^{e-2}(p-1)} \not\equiv 1 \pmod{p^e}.
    \]
    The case $e = 2$ is our choice of $r$. If the assertion holds for $e$, then Euler's theorem modulo $p^{e-1}$ gives
    \[
        r^{p^{e-2}(p-1)} = 1 + cp^{e-1}
    \]
    for some integer $c$ not divisible by $p$. Raising both sides to the $p$th power gives
    \[
        r^{p^{e-1}(p-1)} = (1 + cp^{e-1})^p \equiv 1 + cp^e \not\equiv 1 \pmod{p^{e+1}},
    \]
    proving the induction step. Now let $d$ be the order of $r$ modulo $p^e$. Then $d \mid \varphi(p^e) = p^{e-1}(p-1)$, and reducing modulo $p$ shows $p-1 \mid d$. Hence $d = p^j(p-1)$ for some $0 \leq j \leq e-1$. If $j < e-1$, then $d \mid p^{e-2}(p-1)$, contradicting the congruence just proved. Thus $d = p^{e-1}(p-1) = \varphi(p^e)$, so $r$ is primitive modulo $p^e$. \\
    Finally, let $r$ be primitive modulo $p^e$. Replacing $r$ by $r+p^e$ if necessary, we may assume that $r$ is odd; this does not change its residue modulo $p^e$. Let $d$ be the order of $r$ modulo $2p^e$. Since $r^d \equiv 1 \pmod{2p^e}$ implies $r^d \equiv 1 \pmod{p^e}$, we have $\varphi(p^e) \mid d$. Also $d \mid \varphi(2p^e) = \varphi(p^e)$. Therefore $d = \varphi(2p^e)$, and $r$ is primitive modulo $2p^e$. \\
    We have proved that primitive roots exist exactly for $n = 2$, $4$, $p^e$, and $2p^e$ with $p$ odd. Therefore $\operatorname{Aut}(\Z_n)$ is cyclic exactly for the same values of $n$.
\end{proof}

\noindent The number-theoretic primitive-root classification used in this proof follows [9, Section 8.3].

\subsubsection*{Fixed and Characteristic Subgroups}
\begin{itemize}
    \item Let $\varphi \in \operatorname{Aut}(G)$ and $H \leq G$. We say $H$ is \emph{fixed by} $\varphi$ if $\varphi(H) = H$; in this case the restriction is an automorphism of $H$.
    \item If $L \leq \operatorname{Aut}(G)$, $H$ is \emph{fixed by} $L$ if $H$ is fixed by every $\varphi \in L$.
    \item $H$ is normal in $G$ iff $H$ is fixed by $\operatorname{Inn}(G)$.
    \item $H$ is a \emph{characteristic subgroup} of $G$ (or $H$ \emph{characteristic} in $G$) if $H$ is fixed by $\operatorname{Aut}(G)$. (Some authors write $H \operatorname{char} G$.)
    \item Example: $Z(G)$ is always characteristic. For if $x \in Z(G)$ and $\varphi \in \operatorname{Aut}(G)$, then $\varphi(x)y = \varphi(x \varphi^{-1}(y)) = \varphi(\varphi^{-1}(y)x) = y\varphi(x)$ for any $y \in G$, showing $\varphi(x) \in Z(G)$.
    \item Characteristic subgroups are normal, but the converse fails. In fact, an infinite abelian group need not have any non-trivial proper characteristic subgroups (Exercise~\ref{ex:2.5}).
\end{itemize}

\begin{lemma}
\label{lem:2.4}
If $K$ is a characteristic subgroup of $H$ and $H$ is a characteristic subgroup of $G$, then $K$ is a characteristic subgroup of $G$.
\end{lemma}
\begin{proof}
    Let $\varphi \in \operatorname{Aut}(G)$. We need to show that $\varphi(K) = K$. Since $H$ is characteristic in $G$, we have $\varphi(H) = H$. Hence the restriction $\varphi|_H \colon H \to H$ is an automorphism of $H$: it is still a homomorphism, it is injective because $\varphi$ is injective, and it is surjective onto $H$ because $\varphi(H) = H$. \\
    Now $K$ is characteristic in $H$, so every automorphism of $H$ fixes $K$. Applying this to $\varphi|_H$, we obtain $(\varphi|_H)(K) = K$. But $K \leq H$, so $(\varphi|_H)(K) = \varphi(K)$. Therefore $\varphi(K) = K$. Since $\varphi \in \operatorname{Aut}(G)$ was arbitrary, $K$ is characteristic in $G$.
\end{proof}

\subsubsection*{The Derived Group}
\begin{itemize}
    \item Normality is not transitive: if $N \trianglelefteq G$, the restriction to $N$ of an element of $\operatorname{Inn}(G)$ lies in $\operatorname{Aut}(N)$ but need not lie in $\operatorname{Inn}(N)$.
    \item For $x,y \in G$ the \emph{commutator} is $[x,y] = xyx^{-1}y^{-1}$.
    \item The \emph{derived group} $G'$ of $G$ is the subgroup generated by all commutators: $G' = \langle \{[x,y] \mid x,y \in G\} \rangle$.
    \item $G$ is abelian iff $G' = 1$.
    \item If $H \leq G$, then $H' \leq G'$.
    \item In general, $G'$ contains more than just the commutators of elements of $G$.
    \item Since $[x,y]^{-1} = (xyx^{-1}y^{-1})^{-1} = yxy^{-1}x^{-1} = [y,x]$, by Proposition~\ref{prop:1.2} an arbitrary element of $G'$ is a finite product of commutators of elements of $G$.
\end{itemize}

\begin{lemma}
\label{lem:2.5}
Let $G$ be a group. Then $G'$ is characteristic in $G$.
\end{lemma}
\begin{proof}
    Let $\varphi \in \operatorname{Aut}(G)$. For any $x,y \in G$, we have
    \[
        \varphi([x,y])
        = \varphi(xyx^{-1}y^{-1})
        = \varphi(x)\varphi(y)\varphi(x)^{-1}\varphi(y)^{-1}
        = [\varphi(x),\varphi(y)].
    \]
    Thus $\varphi$ sends each commutator of elements of $G$ to another commutator of elements of $G$. Now let $g \in G'$. By Proposition~\ref{prop:1.2} and the preceding observation about inverses of commutators, we may write
    \[
        g = [x_1,y_1]\cdots [x_r,y_r]
    \]
    for some $x_i,y_i \in G$. Applying $\varphi$, we obtain
    \[
        \varphi(g)
        = [\varphi(x_1),\varphi(y_1)]\cdots [\varphi(x_r),\varphi(y_r)] \in G'.
    \]
    Hence $\varphi(G') \leq G'$. Applying the same argument to $\varphi^{-1}$ gives $\varphi^{-1}(G') \leq G'$, and therefore $G' \leq \varphi(G')$. Thus $\varphi(G') = G'$. Since $\varphi$ was arbitrary, $G'$ is characteristic in $G$.
\end{proof}

\begin{proposition}
\label{prop:2.6}
Let $G$ be a group, and let $N \trianglelefteq G$. Then $G/N$ is abelian iff $G' \leq N$.
\end{proposition}
\begin{proof}
    By Lemma~\ref{lem:2.5}, $G'$ is characteristic in $G$, and hence $G' \trianglelefteq G$. We first observe that $G/G'$ is abelian. Indeed, for any $x,y \in G$,
    \[
        [xG',yG'] = [x,y]G' = G',
    \]
    since $[x,y] \in G'$. Thus every commutator in $G/G'$ is trivial, so $G/G'$ is abelian. \\
    Suppose first that $G' \leq N$. Since $G' \leq N$, the set $N/G'$ is a subgroup of $G/G'$. Moreover, it is normal in $G/G'$: if $g \in G$ and $n \in N$, then
    \[
        (gG')(nG')(gG')^{-1} = gng^{-1}G' \in N/G',
    \]
    because $N \trianglelefteq G$. Thus $N/G' \trianglelefteq G/G'$, and by \nameref{thm:second-iso},
    \[
        G/N \cong (G/G')/(N/G').
    \]
    Since $G/G'$ is abelian, every quotient of $G/G'$ is abelian. Hence $G/N$ is abelian. \\
    Conversely, suppose that $G/N$ is abelian. For any $x,y \in G$, the commutator of $xN$ and $yN$ is trivial, so
    \[
        N = [xN,yN] = [x,y]N.
    \]
    Therefore $[x,y] \in N$ for all $x,y \in G$. Since $N$ contains every commutator of elements of $G$, it contains the subgroup generated by all such commutators. Hence $G' \leq N$.
\end{proof}

\subsubsection*{External Direct Products}
\begin{itemize}
    \item Let $n \in \N$ and let $G_1, \dots, G_n$ be groups. The Cartesian product $G_1 \times \cdots \times G_n$ becomes a group under componentwise multiplication $(g_1,\dots,g_n)(g_1',\dots,g_n') = (g_1 g_1', \dots, g_n g_n')$. The identity is $(1,\dots,1)$ and inverses are $(g_1^{-1},\dots,g_n^{-1})$.
    \item This is the (\emph{external}) \emph{direct product} of $G_1,\dots,G_n$.
    \item Order of factors is irrelevant: $G_1 \times \cdots \times G_n \cong G_{\rho(1)} \times \cdots \times G_{\rho(n)}$ for any $\rho \in S_n$.
    \item For $G = G_1 \times \cdots \times G_n$:
    \begin{itemize}
        \item Each $H_i = \{(1,\dots,g_i,\dots,1) \mid g_i \in G_i\}$ is a normal subgroup of $G$ isomorphic with $G_i$. Moreover $G/H_i$ is isomorphic with the direct product of the remaining $G_j$.
        \item Every $g \in G$ has a unique expression $g = h_1 \cdots h_n$ with $h_i \in H_i$; if $g = (g_1,\dots,g_n)$ then $h_i = (1,\dots,g_i,\dots,1)$.
        \item If the $G_i$ are finite, then $|G| = |G_1| \cdots |G_n|$.
    \end{itemize}
\end{itemize}

\subsubsection*{Internal Direct Products}
\begin{itemize}
    \item Suppose $G$ has subgroups $H_1, \dots, H_n$ such that:
    \begin{enumerate}
        \item $H_i \trianglelefteq G$ for each $1 \leq i \leq n$.
        \item Every $g \in G$ has a unique expression $g = h_1 \cdots h_n$ where $h_i \in H_i$.
    \end{enumerate}
    Conditions (1) and (2) imply:
    \begin{enumerate}
        \setcounter{enumi}{2}
        \item $G = H_1 \cdots H_n$.
        \item $H_i \cap H_1 \cdots H_{i-1} H_{i+1} \cdots H_n = 1$ for each $i$.
        \item If $i \neq j$, then elements of $H_i$ commute with elements of $H_j$.
        \item If $g = h_1 \cdots h_n$ and $g' = h_1' \cdots h_n'$ with $h_i, h_i' \in H_i$, then $gg' = (h_1 h_1') \cdots (h_n h_n')$.
    \end{enumerate}
    \item Under these conditions, there is a unique isomorphism from $G$ to the external direct product $H_1 \times \cdots \times H_n$ sending $H_i$ to $1 \times \cdots \times H_i \times \cdots \times 1$. Call $G$ the (\emph{internal}) \emph{direct product} of $H_1,\dots,H_n$ and write $G = H_1 \times \cdots \times H_n$ (slight abuse of notation).
    \item If (1) holds, then (2) holds iff both (3) and (4) hold; so to check a direct product, it suffices to verify either (1) and (2) or (1), (3), and (4). (For $n = 2$, (4) reduces to $H_1 \cap H_2 = 1$.)
\end{itemize}

\begin{lemma}
\label{lem:2.7}
Let $G$ be a group having normal subgroups $H$ and $K$ such that $G = HK$. Then $G/H \cap K \cong H/H \cap K \times K/H \cap K$.
\end{lemma}
\begin{proof}
    Put $M = H \cap K$. By Proposition~\ref{prop:1.7}, we have $M \trianglelefteq G$, and hence $G/M$, $H/M$, and $K/M$ are quotient groups. By \nameref{thm:correspondence}, the subgroups $H/M$ and $K/M$ are normal in $G/M$. \\
    We first compute their intersection. Suppose that $xM \in (H/M) \cap (K/M)$. Since $xM \in H/M$, there is some $h \in H$ such that $xM = hM$, and because $M \leq H$ this implies $x \in H$. Similarly, $x \in K$. Thus $x \in H \cap K = M$, and so $xM = M$. Therefore
    \[
        (H/M) \cap (K/M) = 1.
    \]
    Next let $gM \in G/M$. Since $G = HK$, we may write $g = hk$ with $h \in H$ and $k \in K$. Then
    \[
        gM = hkM = (hM)(kM),
    \]
    so $(G/M) = (H/M)(K/M)$. Hence $G/M$ is the internal direct product of $H/M$ and $K/M$, and therefore is isomorphic with the corresponding external direct product. Thus
    \[
        G/(H \cap K) \cong H/(H \cap K) \times K/(H \cap K).
    \]
\end{proof}

\begin{lemma}
\label{lem:2.8}
Let $n \in \N$, and write $n = p_1^{a_1} \cdots p_r^{a_r}$ where the $p_i$ are distinct primes and the $a_i$ are positive integers. Then we have $\Z_n \cong \Z_{p_1^{a_1}} \times \cdots \times \Z_{p_r^{a_r}}$.
\end{lemma}
\begin{proof}
    If $n = 1$, there is nothing to prove, so assume $n > 1$. For each $i$, let
    \[
        P_i = \langle x_i \rangle \cong \Z_{p_i^{a_i}},
    \]
    and put $P = P_1 \times \cdots \times P_r$. Then
    \[
        |P| = |P_1|\cdots |P_r| = p_1^{a_1}\cdots p_r^{a_r} = n.
    \]
    Let $x = (x_1,\dots,x_r) \in P$. For $k \in \N$, we have
    \[
        x^k = (1,\dots,1)
    \]
    iff $x_i^k = 1$ for every $i$, iff $p_i^{a_i} \mid k$ for every $i$. Since the prime powers $p_i^{a_i}$ are pairwise coprime, this is equivalent to $n \mid k$. Hence $o(x) = n$. \\
    Therefore $\langle x \rangle$ has $n$ elements and lies in $P$, which also has $n$ elements. Thus $\langle x \rangle = P$, so $P$ is cyclic of order $n$. It follows that $P \cong \Z_n$, and substituting the groups $P_i \cong \Z_{p_i^{a_i}}$ gives
    \[
        \Z_n \cong \Z_{p_1^{a_1}} \times \cdots \times \Z_{p_r^{a_r}}.
    \]
\end{proof}

\begin{corollary}
\label{cor:2.9}
If $\gcd(a,b) = 1$, then $\Z_{ab} \cong \Z_a \times \Z_b$.
\end{corollary}
\begin{proof}
    If $a = 1$ or $b = 1$, the assertion is immediate because $\Z_1$ is the trivial group. Thus assume $a,b > 1$. Write
    \[
        a = p_1^{\alpha_1}\cdots p_s^{\alpha_s},
        \qquad
        b = q_1^{\beta_1}\cdots q_t^{\beta_t}
    \]
    as prime-power factorizations. Since $\gcd(a,b) = 1$, the primes $p_1,\dots,p_s$ are distinct from the primes $q_1,\dots,q_t$. Hence
    \[
        ab = p_1^{\alpha_1}\cdots p_s^{\alpha_s}q_1^{\beta_1}\cdots q_t^{\beta_t}
    \]
    is the prime-power factorization of $ab$. Applying Lemma~\ref{lem:2.8} to $ab$ gives
    \[
        \Z_{ab}
        \cong \Z_{p_1^{\alpha_1}} \times \cdots \times \Z_{p_s^{\alpha_s}} \times \Z_{q_1^{\beta_1}} \times \cdots \times \Z_{q_t^{\beta_t}}.
    \]
    Applying Lemma~\ref{lem:2.8} separately to $a$ and to $b$, and regrouping the factors, gives
    \[
        \Z_{ab} \cong \Z_a \times \Z_b.
    \]
\end{proof}

\begin{proposition}
\label{prop:2.10}
Suppose that a finite group $G$ is the direct product of its subgroups $H_1,\dots,H_n$, where the orders $|H_i|$ are pairwise coprime. Then any subgroup $L$ of $G$ is the direct product of $L \cap H_1, \dots, L \cap H_n$.
\end{proposition}
\begin{proof}
    We first prove the case $n = 2$. Write $G = H \times K$, where $\gcd(|H|,|K|) = 1$, and let $L \leq G$. Put
    \[
        A = L \cap H,
        \qquad
        B = L \cap K.
    \]
    Since $H,K \trianglelefteq G$, we have $A,B \trianglelefteq L$. Also $A \cap B \leq H \cap K = 1$, so $A \cap B = 1$. It remains to show that $L = AB$. \\
    Let $g \in L$. Since $G = H \times K$, we may write $g = hk$ with $h \in H$ and $k \in K$, and $h$ and $k$ commute. By \nameref{thm:lagrange}, $o(h) \mid |H|$ and $o(k) \mid |K|$, so $o(h)$ and $o(k)$ are coprime. Put $a = o(h)$ and $b = o(k)$. If $(hk)^m = 1$, then, since $h$ and $k$ commute,
    \[
        h^m = k^{-m} \in H \cap K = 1.
    \]
    Hence $a \mid m$ and $b \mid m$, and because $\gcd(a,b) = 1$, we have $ab \mid m$. Conversely, $(hk)^{ab} = 1$, so $o(hk) = ab$. \\
    Now $\langle hk \rangle \leq \langle h \rangle\langle k \rangle$. The subgroup $\langle h \rangle\langle k \rangle$ is the internal direct product of $\langle h \rangle$ and $\langle k \rangle$, so it has order $ab$. Since $\langle hk \rangle$ also has order $ab$, we have
    \[
        \langle hk \rangle = \langle h \rangle\langle k \rangle.
    \]
    Thus $h,k \in \langle hk \rangle = \langle g \rangle \leq L$. Therefore $h \in A$ and $k \in B$, and so $g = hk \in AB$. This proves $L \leq AB$, while $AB \leq L$ is immediate. Hence $L = AB$. Since $A$ and $B$ are normal in $L$, have trivial intersection, and commute with one another, $L = A \times B$. \\
    We now argue by induction on $n$. The case $n = 1$ is immediate, and the case $n = 2$ has just been proved. Suppose $n > 2$ and that the result is known for direct products with $n-1$ factors. Let
    \[
        M = H_1 \times \cdots \times H_{n-1}.
    \]
    Then $G = M \times H_n$, and pairwise coprimeness of the orders $|H_i|$ gives $\gcd(|M|,|H_n|) = 1$. Applying the $n = 2$ case to $G = M \times H_n$ gives
    \[
        L = (L \cap M) \times (L \cap H_n).
    \]
    Inside $M$, the subgroup $L \cap M$ is a subgroup of the direct product $H_1 \times \cdots \times H_{n-1}$, whose factor orders are still pairwise coprime. By the induction hypothesis,
    \[
        L \cap M = (L \cap H_1) \times \cdots \times (L \cap H_{n-1}).
    \]
    Substituting this into the preceding display gives
    \[
        L = (L \cap H_1) \times \cdots \times (L \cap H_n),
    \]
    as required.
\end{proof}

\subsubsection*{Internal Semidirect Products}
\begin{itemize}
    \item Suppose $G$ has subgroup $H$ and normal subgroup $N$ with $G = NH$ and $N \cap H = 1$. Then $G$ is the (\emph{internal}) \emph{semidirect product} of $N$ by $H$, written $G = N \rtimes H$. (Notation common but not standard; alternatives include $N \rtimes H$, $H \ltimes N$, no notation.)
    \item If in addition $H \trianglelefteq G$, then $G$ is the direct product of $N$ and $H$.
    \item Example: $G = S_3$, $N = A_3$, $H = \langle (1\,2) \rangle$. Then $G = N \rtimes H$, but $H$ is not normal, so $G$ is not the direct product.
    \item Properties of $G = N \rtimes H$:
    \begin{itemize}
        \item $H \cong H/(N \cap H) \cong NH/N = G/N$ by \nameref{thm:first-iso}; if $G$ is finite, $|G| = |N||G:N| = |N||H|$.
        \item Each $x \in G$ has a unique expression $x = nh$ with $n \in N$, $h \in H$.
        \item If $x = n_1 h_1$, $y = n_2 h_2$, then $xy = n_1 h_1 n_2 h_1^{-1} \cdot h_1 h_2 = n' h'$ with $n' = n_1 h_1 n_2 h_1^{-1} \in N$ and $h' = h_1 h_2 \in H$.
        \item For $h \in H$, conjugation by $h$ maps $N$ to $N$, so $\varphi_h : N \to N$, $\varphi_h(n) = hnh^{-1}$, is an automorphism of $N$, and $\varphi_h \circ \varphi_{h'} = \varphi_{hh'}$. Hence the \emph{conjugation homomorphism} $\varphi : H \to \operatorname{Aut}(N)$, $\varphi(h) = \varphi_h$.
        \item $(n_1 h_1)(n_2 h_2) = n_1 \varphi(h_1)(n_2) \cdot h_1 h_2$, so the multiplication of $G$ is determined by the operations of $N$ and $H$ together with $\varphi$.
        \item If $\varphi$ is trivial, then $hnh^{-1} = n$ for all $n \in N$, $h \in H$, so $H \trianglelefteq G$ and $G = N \times H$. Conversely, if $G = N \times H$, then $H$ commutes with $N$, so $\varphi$ is trivial.
        \item If $\varphi$ is non-trivial, $G$ is non-abelian: there exist $h \in H$, $n \in N$ with $hnh^{-1} = \varphi(h)(n) \neq n$.
    \end{itemize}
\end{itemize}

\subsubsection*{External Semidirect Products}
\begin{itemize}
    \item Given groups $N$, $H$ and a homomorphism $\varphi : H \to \operatorname{Aut}(N)$, define a binary operation on $N \times H$ by
    \[
    (n_1, h_1)(n_2, h_2) = (n_1 \varphi(h_1)(n_2), h_1 h_2).
    \]
    This makes $N \times H$ a group with identity $(1,1)$ and inverse $(n,h)^{-1} = (\varphi(h^{-1})(n^{-1}), h^{-1})$.
    \item Call this the (\emph{external}) \emph{semidirect product} of $N$ by $H$ corresponding to $\varphi$, denoted $G = N \rtimes_{\varphi} H$.
    \item This generally differs from the direct product unless $\varphi$ is trivial.
    \item $\mathcal{N} = N \times 1$ and $\mathcal{H} = 1 \times H$ in $G$ are isomorphic copies of $N$, $H$ respectively. Computation shows $\mathcal{N} \trianglelefteq G$, $G = \mathcal{NH}$, $\mathcal{N} \cap \mathcal{H} = 1$, so $G$ is the internal semidirect product of $\mathcal{N}$ by $\mathcal{H}$.
    \item Identifying $\mathcal{N}$ with $N$ and $\mathcal{H}$ with $H$, the conjugation homomorphism from $\mathcal{H}$ to $\operatorname{Aut}(\mathcal{N})$ corresponds to the original $\varphi$.
    \item Write $N \rtimes_{\varphi} H = \{nh \mid n \in N, h \in H\}$ with multiplication $(n_1 h_1)(n_2 h_2) = n_1 \varphi(h_1)(n_2) \cdot h_1 h_2$. Note $hnh^{-1} = \varphi(h)(n)$.
    \item Distinct $\varphi, \psi : H \to \operatorname{Aut}(N)$ need not yield isomorphic semidirect products. The next two propositions are partial converses.
\end{itemize}

\begin{proposition}
\label{prop:2.11}
Let $H$ be a cyclic group and let $N$ be an arbitrary group. If $\varphi$ and $\psi$ are monomorphisms from $H$ to $\operatorname{Aut}(N)$ such that $\varphi(H) = \psi(H)$, then we have $N \rtimes_{\varphi} H \cong N \rtimes_{\psi} H$.
\end{proposition}
\begin{proof}
    Write $H = \langle x \rangle$. Since $\varphi(H) = \psi(H)$, the automorphisms $\varphi(x)$ and $\psi(x)$ generate the same cyclic subgroup of $\operatorname{Aut}(N)$. Hence there exist $a,b \in \Z$ such that
    \[
        \varphi(x)^a = \psi(x),
        \qquad
        \psi(x)^b = \varphi(x).
    \]
    It follows that, for every $h \in H$,
    \[
        \varphi(h^a) = \psi(h),
        \qquad
        \psi(h^b) = \varphi(h).
    \]
    Indeed, if $h = x^m$, then $\varphi(h^a) = \varphi(x^{am}) = \varphi(x)^{am} = \psi(x)^m = \psi(h)$, and the second equality is similar. \\
    Define
    \[
        \tau \colon N \rtimes_{\psi} H \to N \rtimes_{\varphi} H,
        \qquad
        \tau(nh) = nh^a.
    \]
    We show that $\tau$ is a homomorphism. Let $n_1,n_2 \in N$ and $h_1,h_2 \in H$. Multiplication in $N \rtimes_{\psi} H$ gives
    \[
        (n_1h_1)(n_2h_2) = n_1\psi(h_1)(n_2)h_1h_2.
    \]
    Since $H$ is cyclic, it is abelian, so $(h_1h_2)^a = h_1^a h_2^a$. Therefore
    \[
        \tau((n_1h_1)(n_2h_2))
        = n_1\psi(h_1)(n_2)(h_1h_2)^a
        = n_1\varphi(h_1^a)(n_2)h_1^a h_2^a.
    \]
    But this is exactly the product of $n_1h_1^a$ and $n_2h_2^a$ in $N \rtimes_{\varphi} H$, so
    \[
        \tau((n_1h_1)(n_2h_2)) = \tau(n_1h_1)\tau(n_2h_2).
    \]
    Thus $\tau$ is a homomorphism. Similarly, the map
    \[
        \lambda \colon N \rtimes_{\varphi} H \to N \rtimes_{\psi} H,
        \qquad
        \lambda(nh) = nh^b
    \]
    is a homomorphism. \\
    It remains to check that these two homomorphisms are inverse to one another. Since
    \[
        \varphi(x) = \psi(x)^b = (\varphi(x)^a)^b = \varphi(x^{ab}),
    \]
    and $\varphi$ is injective, we have $x = x^{ab}$. Hence $h = h^{ab}$ for every $h \in H$. Similarly, using the injectivity of $\psi$, we get $h = h^{ba}$ for every $h \in H$. Therefore
    \[
        (\tau \circ \lambda)(nh) = \tau(nh^b) = nh^{ba} = nh,
        \qquad
        (\lambda \circ \tau)(nh) = \lambda(nh^a) = nh^{ab} = nh.
    \]
    Thus $\tau$ and $\lambda$ are mutually inverse isomorphisms, and so $N \rtimes_{\varphi} H \cong N \rtimes_{\psi} H$.
\end{proof}

\begin{proposition}
\label{prop:2.12}
Let $N$ and $H$ be groups, let $\psi : H \to \operatorname{Aut}(N)$ be a homomorphism, and let $f \in \operatorname{Aut}(N)$. If $\hat{f}$ is the inner automorphism of $\operatorname{Aut}(N)$ induced by $f$, then $N \rtimes_{\hat{f} \circ \psi} H \cong N \rtimes_{\psi} H$.
\end{proposition}
\begin{proof}
    By definition, the inner automorphism $\hat{f}$ of $\operatorname{Aut}(N)$ sends $\alpha \in \operatorname{Aut}(N)$ to $f \circ \alpha \circ f^{-1}$. Define
    \[
        \theta \colon N \rtimes_{\psi} H \to N \rtimes_{\hat{f} \circ \psi} H,
        \qquad
        \theta(nh) = f(n)h.
    \]
    We show that $\theta$ is a homomorphism. Let $n_1,n_2 \in N$ and $h_1,h_2 \in H$. Multiplication in $N \rtimes_{\psi} H$ gives
    \[
        (n_1h_1)(n_2h_2) = n_1\psi(h_1)(n_2)h_1h_2.
    \]
    Hence
    \[
        \theta((n_1h_1)(n_2h_2))
        = f(n_1\psi(h_1)(n_2))h_1h_2
        = f(n_1)f(\psi(h_1)(n_2))h_1h_2.
    \]
    On the other hand, multiplication in $N \rtimes_{\hat{f} \circ \psi} H$ gives
    \[
        \theta(n_1h_1)\theta(n_2h_2)
        = (f(n_1)h_1)(f(n_2)h_2)
        = f(n_1)(\hat{f} \circ \psi)(h_1)(f(n_2))h_1h_2.
    \]
    Since
    \[
        (\hat{f} \circ \psi)(h_1)(f(n_2))
        = (f \circ \psi(h_1) \circ f^{-1})(f(n_2))
        = f(\psi(h_1)(n_2)),
    \]
    the two expressions are equal. Thus $\theta$ is a homomorphism. \\
    Finally, the map
    \[
        N \rtimes_{\hat{f} \circ \psi} H \to N \rtimes_{\psi} H,
        \qquad
        nh \mapsto f^{-1}(n)h
    \]
    is inverse to $\theta$. Therefore $\theta$ is an isomorphism, and so $N \rtimes_{\hat{f} \circ \psi} H \cong N \rtimes_{\psi} H$.
\end{proof}

\subsubsection*{Dihedral Groups}
\begin{itemize}
    \item Let $N \cong \Z_n$ ($n \in \N$), $H \cong \Z_2$, $\varphi : H \to \operatorname{Aut}(N)$ send the generator of $H$ to the automorphism sending each element of $N$ to its inverse (so $\varphi(h) = \sigma_{n-1}$).
    \item The group $N \rtimes_{\varphi} H$ is the \emph{dihedral group} of order $2n$, denoted $D_{2n}$. (Some authors write $D_n$.)
    \item $D_{2n}$ is non-abelian iff $n > 2$; when $n = 2$, $\varphi$ is trivial and $D_4 \cong \Z_2 \times \Z_2$.
    \item Geometrically $D_{2n}$ is the symmetry group of a regular $n$-gon: the generator of $N$ is rotation by $2\pi/n$, the generator of $H$ is reflection through a fixed axis.
    \item Also: the \emph{infinite dihedral group} $D_\infty = N \rtimes_{\varphi} H$ where $N \cong \Z$ and $H, \varphi$ as above.
\end{itemize}

\begin{proposition}
\label{prop:2.13}
Let $s$ and $t$ be elements of order $2$ in a group $G$. (Such elements are called \emph{involutions}.) Then $\langle s,t \rangle$ is a dihedral group; in particular, $\langle s,t \rangle = \langle st \rangle \rtimes \langle s \rangle$.
\end{proposition}
\begin{proof}
    Put
    \[
        L = \langle s,t \rangle,
        \qquad
        r = st,
        \qquad
        N = \langle r \rangle,
        \qquad
        H = \langle s \rangle.
    \]
    Since $s$ and $t$ have order $2$, we have $s^{-1} = s$ and $t^{-1} = t$. Therefore
    \[
        srs^{-1} = s(st)s = ts = (st)^{-1} = r^{-1}.
    \]
    Similarly,
    \[
        trt^{-1} = t(st)t = ts = r^{-1}.
    \]
    Thus both generators $s$ and $t$ of $L$ conjugate $r$ to $r^{-1}$. Hence both generators normalize $N = \langle r \rangle$, and so
    \[
        N \trianglelefteq L.
    \]
    Also $H = \langle s \rangle \leq L$. We next show that $L = NH$. Since $r = st \in N$ and $s \in H$, we have
    \[
        t = sr \in HN.
    \]
    But $N \trianglelefteq L$, so $HN = NH$. Hence both $s$ and $t$ lie in $NH$, and therefore
    \[
        L = \langle s,t \rangle \leq NH.
    \]
    The reverse inclusion is immediate from $N,H \leq L$, so $L = NH$. \\
    It remains to check that $N \cap H = 1$. Since $H = \langle s \rangle$ and $s$ has order $2$, the only possible non-identity element of $H$ is $s$. Suppose, for contradiction, that $s \in N$. Then $s$ is a power of $r$. Since $t = sr$, it follows that $t \in N$ as well. Hence $L = \langle s,t \rangle \leq N$, so $L = N$ is cyclic. But a cyclic group has at most one element of order $2$, so $s = t$. Then $r = st = 1$, and hence $N = \langle r \rangle = 1$, contradicting $s \in N$. Therefore $s \notin N$, and so $N \cap H = 1$. \\
    We have shown that
    \[
        L = NH,
        \qquad
        N \trianglelefteq L,
        \qquad
        N \cap H = 1.
    \]
    Hence $L$ is the internal semidirect product $N \rtimes H$. Finally, the conjugation action of $H$ on $N$ sends the generator $r$ of $N$ to $r^{-1}$, so it sends every element of $N$ to its inverse. This is precisely the action used in the definition of a dihedral group. Thus $\langle s,t \rangle$ is dihedral; more explicitly, if $r = st$ has finite order $n$, then $L \cong D_{2n}$, while if $r$ has infinite order, then $L \cong D_\infty$.
\end{proof}

\noindent\textbf{Reference note.}
Reference [7] is Aschbacher's \emph{Finite Group Theory}; the local reference copy is the second edition, and Section 45 is titled ``Involutions in finite groups.'' The point of the reference is that the small fact proved above becomes a major tool once one studies finite simple groups. In a finite group, two involutions $x$ and $y$ generate a dihedral group whose rotation subgroup is $\langle xy \rangle$. This means that the product $xy$ controls the subgroup $\langle x,y \rangle$, and the elementary geometry of dihedral groups gives information about conjugacy among involutions. \\
Aschbacher begins Section 45 by using this exact observation. If $D = \langle x,y \rangle$ and $n = o(xy)$, then $D$ has order $2n$. When $n$ is odd, all involutions of $D$ are conjugate in $D$, so $x$ and $y$ are already conjugate inside the small subgroup they generate. When $n$ is even, there is a unique involution in the rotation subgroup $\langle xy \rangle$, and this central involution gives an additional conjugacy possibility. Thus the parity of $o(xy)$ tells us how involutions can fuse, meaning how they may become conjugate in the larger group. \\
This is important because involutions are unavoidable in non-abelian finite simple groups: by the Feit--Thompson Odd Order Theorem, every non-abelian finite simple group has even order, and therefore has involutions. Section 45 uses dihedral subgroups generated by pairs of involutions to count and control such elements. Two central results in that section are the Brauer--Fowler Theorem and the Thompson Order Formula. The Brauer--Fowler Theorem says, roughly, that fixing the centralizer of an involution leaves only finitely many possible finite simple groups. This suggests that finite simple groups can be attacked by studying the centralizers of their involutions.\\
The Thompson Order Formula is more explicit. When a finite group has at least two conjugacy classes of involutions, it expresses the order of the group in terms of centralizers of involutions and the fusion of involutions inside those centralizers. In one form, if $x_1,\dots,x_k$ represent the conjugacy classes of involutions and $n_i$ counts certain ordered pairs of involutions whose product-generated subgroup contains $x_i$, then
\[
    |G| = |C_G(x_1)|\,|C_G(x_2)|\sum_{i=1}^k \frac{n_i}{|C_G(x_i)|}.
\]
The important point is not the formula itself at this stage, but the method behind it: because pairs of involutions generate dihedral groups, their products and conjugacy relations can be counted very precisely. Thus Proposition~\ref{prop:2.13} is not just a cute exercise. It is one of the basic mechanisms behind the local analysis of finite simple groups, where one tries to understand a group by studying centralizers and normalizers of special subgroups.

\subsection*{Exercises}

\begin{exercise}
\label{ex:2.1}
Let $H$ be a subgroup of a cyclic group $G$. Show that every automorphism of $H$ is the restriction to $H$ of an automorphism of $G$.
\end{exercise}
\begin{solution}
    Let $G = \langle g \rangle$ and let $H \leq G$. Since every subgroup of a cyclic group is cyclic, we may write $H = \langle g^k \rangle$ for some integer $k$. \\
    First suppose that $G$ is infinite cyclic. If $H = 1$, then the only automorphism of $H$ is the restriction of the identity automorphism of $G$. Otherwise $H = \langle g^k \rangle$ is infinite cyclic, and every automorphism of $H$ sends the generator $g^k$ either to $g^k$ or to $g^{-k}$. These two automorphisms are respectively the restrictions of the automorphisms of $G$ given by $g \mapsto g$ and $g \mapsto g^{-1}$. Hence the claim holds when $G$ is infinite cyclic. \\
    Now suppose that $G$ is finite of order $n$, and put $\lvert H \rvert = m$. Let $\psi \in \operatorname{Aut}(H)$. Since $g^k$ generates $H$, we have
    \[
        \psi(g^k) = (g^k)^a
    \]
    for some integer $a$ with $\gcd(a,m) = 1$. We want to choose an automorphism $\varphi \in \operatorname{Aut}(G)$ whose restriction to $H$ is $\psi$. By Proposition~\ref{prop:2.1}, an automorphism of $G$ has the form $\varphi(g) = g^b$ for some integer $b$ with $\gcd(b,n) = 1$. For such a map,
    \[
        \varphi(g^k) = \varphi(g)^k = (g^b)^k = g^{bk} = (g^k)^b.
    \]
    Thus $\varphi$ will agree with $\psi$ on $H$ provided $b \equiv a \pmod{m}$, since $g^k$ has order $m$. \\
    It remains to choose $b$ so that $b \equiv a \pmod{m}$ and $\gcd(b,n) = 1$. Let $q$ be the product of the distinct prime divisors of $n$ which do not divide $m$; if there are no such primes, take $q = 1$. Then $\gcd(m,q) = 1$, so by the Chinese Remainder Theorem there is an integer $b$ such that
    \[
        b \equiv a \pmod{m},
        \qquad
        b \equiv 1 \pmod{q}.
    \]
    If $p$ is a prime divisor of $n$ and $p \mid m$, then $p \nmid a$, since $\gcd(a,m) = 1$, and hence $p \nmid b$ because $b \equiv a \pmod{m}$. If $p \mid n$ but $p \nmid m$, then $p \mid q$, so $b \equiv 1 \pmod{p}$, and again $p \nmid b$. Therefore no prime divisor of $n$ divides $b$, so $\gcd(b,n) = 1$. \\
    Define $\varphi(g) = g^b$. Then $\varphi \in \operatorname{Aut}(G)$ by Proposition~\ref{prop:2.1}, and
    \[
        \varphi(g^k) = (g^k)^b = (g^k)^a = \psi(g^k).
    \]
    Since $g^k$ generates $H$, it follows that $\varphi|_H = \psi$. Thus every automorphism of $H$ is the restriction to $H$ of an automorphism of $G$.
\end{solution}

\begin{exercise}
\label{ex:2.2}
Show that $\operatorname{Aut}(\Z_8) \cong \Z_2 \times \Z_2$.
\end{exercise}
\begin{solution}
    By Proposition~\ref{prop:2.1},
    \[
        \operatorname{Aut}(\Z_8) \cong (\Z/8\Z)^{\times}.
    \]
    The units modulo $8$ are precisely the residue classes represented by integers coprime to $8$, so
    \[
        (\Z/8\Z)^{\times} = \{1 + 8\Z, 3 + 8\Z, 5 + 8\Z, 7 + 8\Z\}.
    \]
    Each non-identity element has order $2$, since
    \[
        3^2 \equiv 5^2 \equiv 7^2 \equiv 1 \pmod{8}.
    \]
    Thus $(\Z/8\Z)^{\times}$ has four elements and every non-identity element has order $2$. Therefore it is isomorphic to $\Z_2 \times \Z_2$, and hence
    \[
        \operatorname{Aut}(\Z_8) \cong \Z_2 \times \Z_2.
    \]
\end{solution}

\begin{exercise}
\label{ex:2.3}
Show that $\operatorname{Aut}(\Z_{p^2}) \cong \Z_{p^2-p}$ for $p$ a prime. (\textsc{Hint:} Let $m$ be a primitive root modulo $p$, and show that either $m$ or $m+p$ be a primitive root modulo $p^2$.)
\end{exercise}
\begin{solution}
    If $p = 2$, then Proposition~\ref{prop:2.1} gives
    \[
        \operatorname{Aut}(\Z_4) \cong (\Z/4\Z)^{\times} = \{1 + 4\Z, 3 + 4\Z\} \cong \Z_2,
    \]
    and this is $\Z_{p^2 - p}$. Hence we may assume that $p$ is odd. \\
    By Proposition~\ref{prop:2.1},
    \[
        \operatorname{Aut}(\Z_{p^2}) \cong (\Z/p^2\Z)^{\times}.
    \]
    The latter group has order $\varphi(p^2) = p^2 - p = p(p - 1)$. It remains to show that it has an element of order $p(p - 1)$. By Proposition~\ref{prop:2.2}, there is a primitive root $m$ modulo $p$, so $m$ has order $p - 1$ modulo $p$. Consider $m$ as a unit modulo $p^2$. If $d$ is its order modulo $p^2$, then $d \mid p(p - 1)$ by Euler's theorem. Also, reducing $m^d \equiv 1 \pmod{p^2}$ modulo $p$ gives $m^d \equiv 1 \pmod{p}$, so $p - 1 \mid d$. Therefore $d$ is either $p - 1$ or $p(p - 1)$. If $d = p(p - 1)$, then $m$ is already a primitive root modulo $p^2$, so $(\Z/p^2\Z)^{\times}$ is cyclic of order $p(p - 1)$. \\
    It remains to handle the case $d = p - 1$. Then
    \[
        m^{p - 1} \equiv 1 \pmod{p^2}.
    \]
    We claim that $m + p$ has order $p(p - 1)$ modulo $p^2$. Since $m + p \equiv m \pmod{p}$, the order of $m + p$ modulo $p^2$ is again divisible by $p - 1$ and divides $p(p - 1)$. Thus it is enough to show that $(m + p)^{p - 1} \not\equiv 1 \pmod{p^2}$. By the binomial theorem,
    \[
        (m + p)^{p - 1}
        \equiv m^{p - 1} + (p - 1)pm^{p - 2}
        \pmod{p^2},
    \]
    since all remaining terms contain a factor $p^2$. Using $m^{p - 1} \equiv 1 \pmod{p^2}$, we obtain
    \[
        (m + p)^{p - 1}
        \equiv 1 + p(p - 1)m^{p - 2}
        \pmod{p^2}.
    \]
    Since $p \nmid m$ and $p \nmid (p - 1)$, the integer $(p - 1)m^{p - 2}$ is not divisible by $p$. Hence $p(p - 1)m^{p - 2}$ is divisible by $p$ but not by $p^2$, and therefore
    \[
        (m + p)^{p - 1} \not\equiv 1 \pmod{p^2}.
    \]
    Thus $m + p$ has order $p(p - 1)$ modulo $p^2$. In either case, $(\Z/p^2\Z)^{\times}$ is cyclic of order $p(p - 1) = p^2 - p$. Therefore
    \[
        \operatorname{Aut}(\Z_{p^2}) \cong \Z_{p^2 - p}.
    \]
\end{solution}

\begin{exercise}
\label{ex:2.4}
Show that if $H \trianglelefteq G$, then any characteristic subgroup of $H$ is normal in $G$.
\end{exercise}
\begin{solution}
    Let $N$ be a characteristic subgroup of $H$. We want to show that $N \trianglelefteq G$. Let $g \in G$ be arbitrary, and consider the inner automorphism
    \[
        \varphi \colon G \to G,
        \qquad
        \varphi(x) = gxg^{-1}.
    \]
    Since $H \trianglelefteq G$, we have
    \[
        \varphi(H) = gHg^{-1} = H.
    \]
    Hence the restriction $\varphi|_H \colon H \to H$ is an automorphism of $H$. Since $N$ is characteristic in $H$, every automorphism of $H$ fixes $N$, so
    \[
        \varphi|_H(N) = N.
    \]
    But $\varphi|_H(N) = gNg^{-1}$. Therefore $gNg^{-1} = N$. Since $g \in G$ was arbitrary, $N \trianglelefteq G$.
\end{solution}

\begin{exercise}
\label{ex:2.5}
Let $F$ be a field, and consider $F$ as a group under its additive operation. Show that $F$ has no non-trivial proper characteristic subgroups.
\end{exercise}
\begin{solution}
\end{solution}

\begin{exercise}
\label{ex:2.6}
Verify the claim made on page 19 that if (1) holds, then (2) holds iff (3) and (4) hold.
\end{exercise}
\begin{solution}
\end{solution}

\begin{exercise}
\label{ex:2.7}
Verify the claim made on page 19 that (1) and (2) together imply (3) through (6).
\end{exercise}
\begin{solution}
\end{solution}

\begin{exercise}
\label{ex:2.8}
Let $G_1$ and $G_2$ be groups, and let $H \leq G_1 \times G_2$. Define
\[
P_1 = \{x \in G_1 \mid (x,y) \in H \text{ for some } y \in G_2\},
\]
\[
I_1 = \{x \in G_1 \mid (x,1) \in H\},
\]
and analogously define subsets $P_2$ and $I_2$ of $G_2$.
\begin{enumerate}
    \item[(a)] Show for $i = 1,2$ that $P_i \leq G_i$ and $I_i \trianglelefteq P_i$.
    \item[(b)] Show that there is a unique isomorphism from $P_1/I_1$ to $P_2/I_2$ sending $xI_1$ to $yI_2$, where $y$ is any element of $G_2$ such that $(x,y) \in H$.
    \item[(c)] Prove \emph{Goursat's theorem}: There is a bijective correspondence between subgroups of $G_1 \times G_2$ and triples $(S_1, S_2, \varphi)$, where $S_i$ is a section of $G_i$ ($i = 1,2$) and $\varphi : S_1 \to S_2$ is an isomorphism. (Recall that a section of a group $G$ is a group $L/M$, where $M \trianglelefteq L \leq G$.)
\end{enumerate}
\end{exercise}
\begin{solution}
\end{solution}

\begin{exercise}
\label{ex:2.9}
(cont.) Use Exercise~\ref{ex:2.8} to give a different proof of Proposition~\ref{prop:2.10}.
\end{exercise}
\begin{solution}
\end{solution}

\begin{exercise}
\label{ex:2.10}
Let $G = N \rtimes H$, and suppose that $N \leq K \leq G$. Show that $K = N \rtimes (K \cap H)$.
\end{exercise}
\begin{solution}
\end{solution}

\subsection*{Further Exercises}

\noindent Let $N$ and $H$ be groups. An \emph{extension} of $N$ by $H$ is a group $E$ along with a monomorphism $i : N \to E$ and an epimorphism $\pi : E \to H$ such that $i(N) = \ker \pi$ (so that $N$ imbeds in $E$ as a normal subgroup, with the quotient group being isomorphic with $H$). We shall usually refer to an extension $(E,i,\pi)$ simply by the group $E$; however, the nature of the maps $i$ and $\pi$ are important in distinguishing between extensions. We identify $N$ with its image under $i$, and $H$ with the quotient of $E$ by $N$. As an example, let $\varphi : H \to \operatorname{Aut}(N)$ be a homomorphism; then the semidirect product $N \rtimes_{\varphi} H$ is an extension of $N$ by $H$ in an obvious way, taking $i$ to be the inclusion map sending $n \in N$ to $(n,1)$ and $\pi$ to be the projection map sending $(n,h)$ to $h$.

\begin{exercise}
\label{ex:2.11}
We say that an extension $E$ of $N$ by $H$ is a \emph{split extension} if there is a homomorphism $t : H \to E$ (called a \emph{splitting map} for the extension) such that $\pi \circ t$ is the identity map on $H$, in which case $t(H)$ will be a transversal for $N$ in $E$. Show that $E$ is a split extension iff it is a semidirect product of $N$ by $H$.
\end{exercise}
\begin{solution}
\end{solution}

\begin{exercise}
\label{ex:2.12}
(cont.) Let $Q$ be the quaternion group of order 8. (We can consider $Q$ as the set $\{\pm 1, \pm i, \pm j, \pm k\}$ with multiplication given by the rules $i^2 = j^2 = k^2 = -1$ and $ij = k = -ji$.) Show that $Q$ can be realized as a non-trivial extension in four ways---thrice as an extension of $\Z_4$ by $\Z_2$, and once as an extension of $\Z_2$ by $\Z_2 \times \Z_2$---but that none of these extensions is split. (In other words, $Q$ cannot be written non-trivially as a semidirect product.)
\end{exercise}
\begin{solution}
\end{solution}

\noindent If $E$ is an extension of $N$ by $H$, then we cannot expect to find a homomorphism $t : H \to E$ such that $t(H)$ will be a transversal for $N$ in $E$, for if such a $t$ existed then $E$ would be split. However, since $H \cong E/N$, we can always find a set map $t : H \to E$ whose image is a transversal for $N$; such a map is called a \emph{section} of the extension. Moreover, we can always choose $t$ so that $t(1) = 1$, in which case we say that $t$ is \emph{normalized}. (We use normalized sections instead of arbitrary sections in order to keep the notational complexity to a minimum.)

\begin{exercise}
\label{ex:2.13}
(cont.) Let $t$ be a normalized section of an extension $E$. Let $\Psi : E \to \operatorname{Aut}(E)$ be the homomorphism sending an element of $E$ to the corresponding inner automorphism of $E$. We shall, for $x \in E$, regard $\Psi(x)$ as being an automorphism of $N$, which is possible since $N \trianglelefteq E$. Define set maps $f : H \times H \to N$ and $\varphi : H \to \operatorname{Aut}(N)$ by
\[
f(\alpha, \beta) = t(\alpha) t(\beta) t(\alpha \beta)^{-1},
\]
\[
\varphi(\alpha) = \Psi(t(\alpha)).
\]
We call $(f, \varphi)$ the \emph{factor pair} arising from $t$. Show that $(f, \varphi)$ has the following properties:
\begin{enumerate}
    \item[1] $f(\alpha, 1) = f(1, \alpha) = 1$ for every $\alpha \in H$, and $\varphi(1)$ is the identity in $\operatorname{Aut}(N)$.
    \item[2] $\varphi(\alpha)\varphi(\beta) = \Psi(f(\alpha, \beta))\varphi(\alpha\beta)$ for $\alpha, \beta \in H$.
    \item[3] $f(\alpha, \beta) f(\alpha\beta, \gamma) = \varphi(\alpha)(f(\beta, \gamma)) f(\alpha, \beta\gamma)$ for $\alpha, \beta, \gamma \in H$.
\end{enumerate}
\end{exercise}
\begin{solution}
\end{solution}

\begin{exercise}
\label{ex:2.14}
(cont.) Just as we were able to externalize the notion of semidirect product, so should we be able to externalize the notion of extension; that is, given groups $N$ and $H$ and appropriate additional data, we should be able to construct an extension of $N$ by $H$. Using Exercise~\ref{ex:2.13} as a guide, formulate such an external construction and prove that it works.
\end{exercise}
\begin{solution}
\end{solution}

\noindent We shall return to these ideas in the further exercises to Section 9.

\end{document}
